\begin{frame}{Experimental Setup}
  \begin{itemize}
    \item Initial experiments were conducted with board $6 \times 6$, $5$ mines 
    (density $\approx 13.9\%$) \alert{identical} for DQN and PPO. The metric
    we aim to maximize is the \texttt{win\_rate}, defined as number of episodes
    in which the agent won over the total number of test episodes. In this first phase,
    we were able to achieve a max win rate of $\approx 80.1\%$ for DQN and $\approx 80.2\%$ for
    PPO.
    
    Further experiments were conducted with boards $9 \times 9$ and $10$ mines.
    \item disjoint RNG seed ranges for train / validation / held-out test
    \item final evaluation: $1000$ held-out test episodes per agent
    \item to conduct in a rigorous and ordered manner an experimental performance analysis,
    we used the \href{https://wandb.ai/site/}{Weight and Biases} platform (academic plan) and its Python Library.
    The source code interacting with W\&B to store the complete run logs was produced with
    the help of Claude Opus 4.6 and Gemini 3.5 Flash.
  \end{itemize}
\end{frame}
\begin{frame}{Statistical tests: Fixing DRL algorithm, \texttt{fully\_conv} VS \texttt{global\_skip\_conv}}
Due to the nature of the environment, the metric we aimed to maximize is the win rate.
Statistical test used: McNemar test; it's the equivalent of the
Wilcoxon-signed rank test for binary events. The test focuses only on disagreeing episodes.

\small{$c = $ \texttt{num\_episodes} won by $B$ \hspace{1.5em} $b$ \texttt{num\_episodes} won by $A$.

\texttt{p\_value} = Probability of observing $c$ wins of B over $b+c$ total victories

$ = P(X=c)$, $X \sim \text{Binom}(n=b+c, p=0.5)$. $\mathcal{H}_0: P(A:wins, B:lose) = P(A:lose, B:wins)$}

  \textbf{DQN}
  \vspace{-0.2em}
  {\scriptsize
  \centering
  \setlength{\tabcolsep}{3pt}
  \renewcommand{\arraystretch}{0.85}

  \begin{tabular}{c|cccc}
  \texttt{agent\_seed}
  &
  Win rate (\texttt{fully\_conv})
  &
  Win rate (\texttt{global\_skip})
  &
  $p$-value
  &
  Significant ($\alpha=0.05$)
  \\
  \hline
  43 & 69\% & 63\% & $3.32 \times 10^{-3}$ & yes \\
  44 & 67\% & 67\% & $0.88$                & no  \\
  45 & 68\% & 67\% & $0.60$                & no
  \end{tabular}

\par}
\vspace{0.15em}

\textbf{PPO}
\vspace{-0.2em}
{\scriptsize
\centering
\setlength{\tabcolsep}{3pt}
\renewcommand{\arraystretch}{0.85}

\begin{tabular}{c|cccc}
\texttt{agent\_seed}
&
Win rate (\texttt{fully\_conv})
&
Win rate (\texttt{global\_skip})
&
$p$-value
&
Significant ($\alpha=0.05$)
\\
\hline
43 & 62\% & 72\% & $7.48 \times 10^{-8}$ & yes \\
44 & 65\% & 67\% & $0.28$                & no  \\
45 & 56\% & 65\% & $1.68 \times 10^{-6}$ & yes
\end{tabular}

\par}


\end{frame}
\begin{frame}{Statistical tests: Fixing architecture, DQN vs PPO}
Tabelle, colonna sx con fully conv colonna destra con global skip
\end{frame}
\begin{frame}{Future works}
Se vogliamo aggiungere qualcosa su curriculum learning
e self attention nel ramo globale.
\end{frame}
\begin{frame}{Slide extra}
Extra slide, più che altro se non abbiamo spazio
per i grafici e/o parametri, ma con questa abbiamo già sforato di 1.
\end{frame}